The recursive metric transport proof RMT1--18 also applies here on the
actual moment source and supplies a further propagated refinement.
Its coordinate map is $y\mapsto u=y$, leaving $S=k/2+iu$ unchanged.
Equation (FV.M1) equals its reference measure (RMT2)--(RMT3), including
$c_{1/4}^k=(\sqrt{2\pi})^k$; the arithmetic measure is identically
$w_h^{*k}du$ in both. Therefore all entries of its $M_x$ equal (FV.M2).
Its included relation matrix is $I_NB_N=\widetilde B_N$ in (FV.M3),
and its quotient is this same $J_N$. Substitution proves equality of
both Gram inverses, every $P_N,Q_N$, and the four volumes individually.
In particular its $\mathcal B(x)$ equals the original $F(x)$ in (FV.M5),
with no reversal of an endpoint or alteration of a sign. Its quantities
$E_q,O_q,d_x$ are precisely $S_{\rm AW},S_{\rm SP},d_x$ here.

The measure hypothesis of that result holds literally, rather than being
assigned to an arbitrary coefficient form. Its required positivity
$F(x)>0$ can also be checked in these coordinates: if $F(x)=0$,
the two crossed nested ratios $(q-1,2q)$ and $(q,2q-1)$ both have
determinant ratio one. Their least-norm sections, included in the common
$\mathcal H$ by $I_i,I_j$, then agree, because
$G_i-G_j=(I_iR_i-I_jR_j)^*M_x(I_iR_i-I_jR_j)$ is positive and determinant one
forces equality of the two forms. In the first crossed pair the literal
remainder section contains $1$, which would be orthogonal to
$\chi\mathcal P_q$. The polynomial
$\chi^{\#_k}(S)=\sum_j\overline{a_j}(k-S)^j$ for
$\chi(S)=\sum_j a_jS^j$ equals $\overline{\chi(S)}$ on the same
line. Orthogonality to $\chi\chi^{\#_k}$ would force
$\int|\chi(S)|^2m_x(u)du=0$, contrary to the positive degree-$2q$
moment Gram. This retains the full original relation polynomial.

The further two entries in the current recursive control are calculated
on these same operators. Retain
\[
 \begin{aligned}
 T_q(x)&=\frac{d_x}{2}\|U_x-W_x\|_1,\\
 \Sigma_x&=\operatorname{Tr}(C_x^2)
                -\frac{(\operatorname{Tr}C_x)^2}{2q+1},\\
 S_{\rm HS}(x)&=\sqrt{\Sigma_x
                         \operatorname{Tr}((U_x-W_x)^2)}.
 \end{aligned}\tag{FV.M15a}\label{fvmetric:trace-hs}
\]
The midpoint-centering argument proving (FV.M8) already proves
$|F'(x)|\le T_q(x)\le S_{\rm SP}(x)$ before either trace-norm
majorant is taken. Both inequalities are retained. For the second new
entry, $C_x$ is self-adjoint in $M_x$ because $M_xC_x=M_1-M_0$ is
Hermitian. Its centered operator
$C_x-\operatorname{Tr}(C_x)I/(2q+1)$ has squared Hilbert--Schmidt
norm exactly $\Sigma_x\ge0$, by expanding its square and trace.
Since $\operatorname{Tr}(U_x-W_x)=0$, replacing $C_x$ by this
centered operator leaves the trace pairing $F'(x)$ unchanged.
Cauchy--Schwarz on the finite matrix entries in an $M_x$-orthonormal
basis therefore gives $|F'(x)|\le S_{\rm HS}(x)$, with the other
squared norm exactly the trace in (FV.M8). This basis is used only
to prove the Hilbert-space inequality: both endpoint forms and every
source-coordinate operator remain the original ones. The expressions
in (FV.M15a) equal RMT10a--RMT10b after the coefficient identifications
already proved above.

Put $a=2q-1$. The full-angle estimate of RMT11, whose two complete
crossed angle lists and original zero angle are retained in that proof,
therefore gives the simultaneous quantities
\[
 \begin{aligned}
 A_q(x)&=a\sqrt{1-e^{-F(x)/a}}\,d_x,\\
 \mathfrak J_{h,k}&=\int_0^1
       \min\{S_{\rm AW}(x),T_q(x),S_{\rm SP}(x),S_{\rm HS}(x),A_q(x)\}\,dx
                    \le\mathfrak I_{h,k},\\
 \mathfrak I_{h,k}^{\rm nl}&=\int_0^1\min\Bigl\{d_x,
 \frac{S_{\rm AW}(x)}{a\sqrt{1-e^{-F(x)/a}}},
 \frac{T_q(x)}{a\sqrt{1-e^{-F(x)/a}}},\\
 &\hspace{15mm}
 \frac{S_{\rm SP}(x)}{a\sqrt{1-e^{-F(x)/a}}},
 \frac{S_{\rm HS}(x)}{a\sqrt{1-e^{-F(x)/a}}}\Bigr\}\,dx
                    \le\log\frac{b_{2q+1}}{b_1}.
 \end{aligned}\tag{FV.M15}\label{fvmetric:recursive-integrals}
\]
The exact source identities above identify all five retained estimates
individually with RMT9--RMT13, including the exact trace-norm and
centered Hilbert--Schmidt entries. Each bounds the same $|F'(x)|$, so
their pointwise minimum does also. In particular it is at most the
earlier three-entry minimum and the two-entry $\mathfrak I_{h,k}$;
none of their full proofs is discarded. Strict positivity and compactness
of the fixed path keep the displayed denominators nonzero. RMT13 proves
measurability of the actual range intersection, and (FV.M10) computes
the last integral of $d_x$. Integrating gives
$|\Delta_{h,k}|\le\mathfrak J_{h,k}$ directly.

For completeness the scalar coordinate used to propagate this refinement
is $\mathscr H_a(B)=2\operatorname{arcosh}(e^{B/(2a)})$.
It has derivative $1/(a\sqrt{1-e^{-B/a}})$ and inverse
$2a\log\cosh(z/2)$, by direct differentiation and substitution on the
positive real branches. Apply this derivative to the five bounds in
(FV.M15) and integrate, exactly as in the same-source RMT17, to obtain
\[
 |\mathscr H_a(\mathcal B_{h,k})
           -\mathscr H_a(\mathcal B^\Gamma_{h,k})|
       \le\mathfrak I_{h,k}^{\rm nl}.
 \tag{FV.M16}\label{fvmetric:recursive-coordinate}
\]
Write $H_0=\mathscr H_a(\mathcal B^\Gamma_{h,k})$ and retain the
complete simultaneous endpoints
\[
 \begin{aligned}
 \mathcal B_{h,k}^{\rm lo}
 &=\max\left\{0,\mathcal B^\Gamma_{h,k}-\mathfrak J_{h,k},
 2a\log\cosh\left(\frac{\max\{0,H_0-\mathfrak I_{h,k}^{\rm nl}\}}2\right)
                                                        \right\},\\
 \mathcal B_{h,k}^{\rm hi}
 &=\min\left\{\mathcal B^\Gamma_{h,k}+\mathfrak J_{h,k},
 2a\log\cosh\left(\frac{H_0+\mathfrak I_{h,k}^{\rm nl}}2\right)
                                                        \right\}.
 \end{aligned}\tag{FV.M17}\label{fvmetric:recursive-endpoints}
\]
The inverse is increasing and extends continuously by zero at zero.
Applying it to (FV.M16) and intersecting the resulting interval with
$|\Delta_{h,k}|\le\mathfrak J_{h,k}$ proves
\[
 \mathcal B_{h,k}^{\rm lo}\le\mathcal B_{h,k}
 \le\mathcal B_{h,k}^{\rm hi}
 \le\mathcal B^\Gamma_{h,k}+\mathfrak J_{h,k}
 \le\mathcal B^\Gamma_{h,k}+\mathfrak I_{h,k}.
 \tag{FV.M18}\label{fvmetric:recursive-budget}
\]
These formulas use the complete proof in
Section~\ref{sec:recursive-metric-transport}, with the actual measure
and coefficient correspondences just proved. The preceding generic
signed matrix calculation (FV.M3)--(FV.M10) remains valid on its own
stated domain; the full-angle refinement here is specialized to the
actual original arithmetic/Gamma moment path.
